\documentclass{subfiles}
\begin{document}
    Introduced in 1989 by R. Capocelli, Fibonacci encoding produces self-syncronizing codewords.
        The encoding is based on the Zeckendorf theorem stated below,
        which allows us to represent each positive integer as the sum of Fibonacci numbers.
        Recall that Fibonacci numbers are defined by the recurrence 
        \[
            F_{n} = \begin{cases}
                0, \text{ if } n = 0 \\ 
                1, \text{ if } n = 1 \\ 
                F_{n - 1} + F_{n - 2}, \forall n \ge 2.
            \end{cases}
        \]

    \begin{theorem*}[Zeckendorf's]
        Let \(n\) be an integer. 
            Then, \(n\) is either a Fibonacci number or 
            there exist distinct non-consecutive Fibonacci numbers
            \(F_{i_{1}}, F_{i_{2}, \ldots}\) such that
            \[
                n = F_{i_{1}} + F_{i_{2}} + \cdots 
            \]
            with \(F_{i_{j}} \ge 1\).
    \end{theorem*}
    Zeckendorf's representation has the property that no two consecutive bits 
        are set to 1.

    Let \(n\) be an integer. To encode it using Fibonacci encoding:
        find the greast Fibonacci number lower then or equal to \(n\),
        say it's \(F_{i_{j}}\), and compute the reminder \(r = n - F_{i_{j}}\).
        Starting from the left, set the \(i_{j}\)-th bit to 1.
        Keep computing the reminder replacing at each step \(r\) to \(n\).
        Append a 1 to the sequence.

    Appending a 1 to the representation, allows us to have self-syncronizing codes.
        If errors a occur in trasmission, the affected codeword is discarded 
        without affecting the rest of the bitstream. 

    \begin{example*}
        Consider the integer 177. To begin with compute all Fibonacci numbers 
            lower than or equal to 177, these are: \(F_{1} = 1, F_{2} = 2, F_{3} = 3,
            F_{4} = 5, F_{5} = 8, F_{6} = 13, F_{7} = 21, F_{8} = 34, F_{9} = 55,
            F_{10} = 89, F_{11} = 144\). Take \(F_{11}\) and compute \(r = 177 - 144 = 33\),
            and set the 11-th bit from the left to 1. 
            We have for the moment \(Fib(177) = 00000000001\).
            By computing \(r\) once more we have \(r = 33 - 21 = 12\);
            since the \(7\)-th Fibonacci number was used we update the encoding by setting the 
            \(7\)-th bit to one. Repeating the process we get that \(Fib(177) = 101010100011\).
    \end{example*}

    Decoding is trivial: locate the ``11'' in the codeword, and remove the extra bit.
        Replace each 1 within the codeword with the corresponding Fibonacci number,
        and compute the integer \(x\) as the sum of these values.

    Fibonacci encoding is proved optimal for 
    \[
        p(x) = \frac{1}{x^{1.44}}.
    \]
\end{document}
